\documentclass[11pt]{article}
\usepackage[margin=1in]{geometry}
\usepackage{amsmath,amssymb,amsthm,enumitem,hyperref}
\hypersetup{colorlinks=true,linkcolor=blue,urlcolor=blue,citecolor=blue}
\newtheorem{theorem}{Theorem}
\newcommand{\Q}{\mathcal Q}
\newcommand{\Z}{\mathbb Z}
\newcommand{\supp}{\operatorname{supp}}
\newcommand{\M}{\overline M}
\title{Duality for transmuted preorder $M$-triangles\\\large MathDB 374967}
\author{Principia Math In Progress}
\date{September 18, 2026}
\begin{document}
\maketitle
\begin{abstract}
We prove Conjecture 8.2 of Athanasiadis and Chapoton for every finite
preorder. The proof realizes a rank-and-support enumerator in a
two-variable lattice-point series. The interior points of the associated
cone are a fixed translate of all its lattice points. Stanley reciprocity
then gives exactly the proposed transmuted $M$-triangle duality.
\end{abstract}

\section*{Task 1. Full solution}
Let $\tau$ be a preorder on an $n$-element set $E$. Put
\[
 P_\tau=\{a\in\Z_{\ge0}^E:a(I)\le |I|
                       \text{ for every order ideal $I$ of $\tau$}\},
 \qquad a(I)=\sum_{e\in I}a_e.
\]
Its order is componentwise and its rank is $|a|=a(E)$. Write
$s(a)=|\supp(a)|$. For a finite graded poset $P$, define
\[
 M_P(x,y)=\sum_{a\le b}\mu_P(a,b)x^{|a|}y^{|b|},\qquad
 \M_P(x,y)=M_P\left(\frac{1-y}{1-xy},1-xy\right).
\]
\begin{theorem}\label{thm:main}
For every finite preorder $\tau$ of size $n$,
\[
 \boxed{(xy)^n\M_{P_\tau}(1/x,1/y)=\M_{P_\tau}(y,x).}
\]
No lattice realization of $P_\tau$ is assumed.
\end{theorem}

We may suppose $n\ge1$, since the empty preorder gives both sides
equal to one. Introduce the homogeneous degree-$n$ polynomial
\begin{equation}\label{eq:H}
 H_\tau(X,Y,Z)=\sum_{a\in P_\tau}
                    X^{n-|a|}Y^{s(a)}Z^{|a|-s(a)}.
\end{equation}
All exponents are nonnegative, since $|a|\le n$.
Downward closure makes each interval $[a,b]$ a product of chains.
Its M\"obius function is zero unless every coordinate of $b-a$ is
zero or one, in which case it is $(-1)^{|b|-|a|}$. Summing first
over $a\le b$ gives
\[
 M_{P_\tau}(x,y)=\sum_{b\in P_\tau}(xy)^{|b|}(1-x^{-1})^{s(b)}.
\]
After transmutation this becomes the polynomial identity
\begin{equation}\label{eq:M}
 \M_{P_\tau}(x,y)
 =\sum_{b\in P_\tau}(-y)^{s(b)}(1-x)^{s(b)}(1-y)^{|b|-s(b)}
 =H_\tau(1,y(x-1),1-y).
\end{equation}
We will prove the stronger homogeneous identity
\begin{equation}\label{eq:Hinvolution}
 H_\tau(X,Y,Z)=H_\tau(X+Y-Z,Z,Y).
\end{equation}

\subsection*{A cone and its interior translation}
Set $\sigma=\tau^*$, the reversed preorder. In coordinates
$(t,\alpha,\beta)\in\mathbb R\times\mathbb R^E\times\mathbb R^E$
consider the rational cone
\begin{equation}\label{eq:cone}
 C=\{(t,\alpha,\beta):\alpha,\beta\ge0,\quad
       \beta(I)+\alpha(E\setminus I)\le t
                    \text{ for every ideal $I$ of $\sigma$}\}.
\end{equation}
Here the empty ideal is included. The cases $I=\varnothing,E$
give $|\alpha|,|\beta|\le t$. Thus $C$ is pointed. It is full
dimensional, since
\[
 w=(n+1,\mathbf1,\mathbf1)
\]
strictly satisfies all its inequalities. Every displayed normal
is nonzero, so strict inequalities characterize the interior, even
if some inequalities are redundant. For integer points, they are
$\alpha_e,\beta_e\ge1$ and
$\beta(I)+\alpha(E\setminus I)\le t-1$.
Subtracting $w$ therefore gives the exact equality
\begin{equation}\label{eq:translate}
 C^\circ\cap\Z^{2n+1}=w+(C\cap\Z^{2n+1}).
\end{equation}
Indeed, the left side of each ideal inequality drops by $n$, while
the coordinate $t$ drops by $n+1$.

Let $\alpha_0=t-|\alpha|$. It is a nonnegative integer on lattice
points of $C$. Define the two-variable series
\[
 G(u,v)=\sum_{(t,\alpha,\beta)\in C\cap\Z^{2n+1}}
                                u^{\alpha_0}v^{|\alpha|}.
\]
Every coefficient is finite: fixing $\alpha_0$ and $|\alpha|$
fixes $t$, and both nonnegative vectors have sum at most $t$.
Moreover a nonzero cone vector has positive total degree in this
grading. Stanley's reciprocity theorem for rational cones
\cite[Theorem 1]{BD}, followed by this monomial specialization,
applies without a zero-degree ray. Since $\dim C=2n+1$ and
translation by $w$ changes $(\alpha_0,|\alpha|)$ by $(1,n)$,
\eqref{eq:translate} gives, as rational functions,
\begin{equation}\label{eq:Grec}
 G(u^{-1},v^{-1})=-u v^n G(u,v).
\end{equation}

\subsection*{The numerator of the series}
For fixed nonnegative integers $(\alpha_0,\alpha)$, the possible
$\beta$ form the lattice points of
\begin{equation}\label{eq:section}
 \alpha_0\Delta_E+
       \sum_{e\in E}\alpha_e\Delta_{U_\sigma(e)},\qquad
 \Delta_A=\operatorname{conv}(0,\mathbf e_f:f\in A).
\end{equation}
To check this, the cone inequalities become
$\beta(I)\le\alpha_0+\alpha(I)$ for nonempty ideals.
The transportation criterion for the Minkowski sum is
$\beta(A)\le\alpha_0+\alpha(D_\sigma(A))$ for every nonempty
subset $A$. These conditions are equivalent: take ideals for one
direction and use $A\subseteq D_\sigma(A)$ for the other.

Postnikov's established formula in
\cite[Theorem 2.1 and Lemma 4.1]{AC}, with these arbitrary summand
weights, says that the number of points in \eqref{eq:section} is
\begin{equation}\label{eq:count}
 \sum_{a\in P_\tau}
       \binom{\alpha_0+n-|a|}{n-|a|}
       \prod_{e\in E}\binom{\alpha_e+a_e-1}{a_e}.
\end{equation}
The additional draconian coordinate equals $n-|a|$; its summand
is the universal simplex, which explains the first binomial factor.
For fixed nonnegative $k$, the needed series are
\[
 \sum_{j\ge0}\binom{j+k}{k}u^j=(1-u)^{-k-1},\qquad
 \sum_{j\ge0}\binom{j+k-1}{k}v^j=
 \begin{cases}(1-v)^{-1}&k=0,\\v(1-v)^{-k-1}&k>0.\end{cases}
\]
In the $k=0$ case the binomial coefficient is one, including $j=0$.
Sum \eqref{eq:count} with weights $u^{\alpha_0}v^{|\alpha|}$.
For $r=|a|$ and $s=s(a)$, its contribution is
$(1-u)^{-n+r-1}v^s(1-v)^{-n-r}$. Consequently
\begin{equation}\label{eq:G}
 G(u,v)=\frac{H_\tau(1-v,v(1-u),1-u)}
                    {(1-u)^{n+1}(1-v)^{2n}}.
\end{equation}
The root-polytope interpretation and its one-variable support
enumerator appear in \cite[Section 4.2]{D}; here we retain the
extra grading needed for the $M$-triangle.

\subsection*{Reciprocity becomes transmuted duality}
Homogeneity of $H_\tau$ and \eqref{eq:G} give
\[
 G(u^{-1},v^{-1})=
 -u v^n\frac{H_\tau(u(1-v),1-u,v(1-u))}
                    {(1-u)^{n+1}(1-v)^{2n}}.
\]
Cancel the common factors using \eqref{eq:Grec}. We obtain
\begin{equation}\label{eq:param}
 H_\tau(1-v,v(1-u),1-u)
       =H_\tau(u(1-v),1-u,v(1-u)).
\end{equation}
The first coordinate on the right equals
$(1-v)+v(1-u)-(1-u)$, so this is \eqref{eq:Hinvolution}
on the displayed parametrization. Homogeneity extends it to all
triples. Explicitly, for generic $(X,Y,Z)$ take
\[
 v=Y/Z,\qquad u=1-(Z-Y)/X,\qquad \lambda=XZ/(Z-Y).
\]
Then $(X,Y,Z)=\lambda(1-v,v(1-u),1-u)$. Thus the two polynomials
in \eqref{eq:Hinvolution} agree on a dense open set and agree identically.

Finally, substitute $(X,Y,Z)=(1,x(y-1),1-x)$ into
\eqref{eq:Hinvolution}, and use \eqref{eq:M} and homogeneity:
\begin{align*}
 \M_{P_\tau}(y,x)
 &=H_\tau(1,x(y-1),1-x)\\
 &=H_\tau(xy,1-x,x(y-1))\\
 &=(xy)^n H_\tau\left(1,\frac{1-x}{xy},\frac{y-1}{y}\right)
  =(xy)^n\M_{P_\tau}(1/x,1/y).
\end{align*}
This proves Theorem~\ref{thm:main} as a polynomial identity.

\subsection*{Verification and scope}
The exact verifier checks 78,480 M\"obius intervals and all 389
labeled preorders of sizes 1--4. It checks 2,466 weighted cone sections
containing 74,376 lattice points, including their interior translations;
the source's five-element triangle and larger examples also agree.
All 26 pages of the primary and all 25 pages of the later root-polytope
paper were read. The latter proves support-enumerator duality, a
different statement. Bounded checks found no earlier full target claim
or target journal final. Review is by the authoring agent only.

\section*{Tasks 2--5. Proposed follow-up directions}
\begin{enumerate}[label=\textbf{Task \arabic*.},start=2,leftmargin=*]
\item Describe the coefficient relations in \eqref{eq:Hinvolution}
directly in terms of ranks, support sizes, and slack coordinates.
\item Retain separate weights on the supply coordinates and derive
a multivariate refinement of the transmuted duality.
\item Characterize bipartite graphs whose augmented root cone admits
an interior translation with the same grading shift.
\item Determine which of these reciprocal triangles are realized by
finite graded lattices, addressing the additional existence question.
\end{enumerate}
\begin{thebibliography}{9}
\raggedright
\bibitem{AC} C. A. Athanasiadis and F. Chapoton,
\emph{Polytopes and posets associated to preorders},
\href{https://arxiv.org/abs/2605.26916v1}{arXiv:2605.26916v1}, 2026.
\bibitem{BD} M. Beck and M. Develin,
\emph{On Stanley's reciprocity theorem for rational cones},
\href{https://arxiv.org/abs/math/0409562v3}{arXiv:math/0409562v3}, 2005.
\bibitem{D} Z. Dai, Q. Hou, Z. Liu, W. Thawinrak, and H. Wang,
\emph{Counting Lattice Points in Minkowski Sums of Cross Polytopes},
\href{https://arxiv.org/abs/2608.16037v2}{arXiv:2608.16037v2}, 2026.
\end{thebibliography}
\end{document}
